
The method requires only a pre-trained flow-based model and an ODE/SDE solver. The three samplers share one loop per assimilation step -- drawing $\state^{n}\sim\conddist{\state^{n}}{\obs^{n},\state^{n-1}}$ -- which at each of $M$ pseudo-time steps evaluates velocity and score from the model, forms the observation interpolant $\interpolantobs_\tau$ and likelihood score $\bar S_\tau$ (Lemma~\ref{lem:source_moments}, Eq.~\eqref{eq:interpolant_score_closed}), and integrates the posterior drift $\postdrift_\tau=\drift^{\gdiff}_\tau+\gweight_\tau\bar S_\tau$, adding the Brownian increment when $\gdiff_\tau>0$. The samplers differ only in the source, $\gdiff_\tau$, and $\gweight_\tau$; pseudo-code for the shared loop and the three specializations is in Appendix~\ref{appendix:algorithms} (Algorithms~\ref{alg:unified}--\ref{alg:fm_ode}).

Several points are worth noting. \emph{Cost.} The loop solves no inner integral and backpropagates through no trajectory. With Corollary~\ref{cor:cheap_drift}, $\bar\Sigma_\tau\approx\beta_\tau^2\obscov+\sigpath_\tau^2\obsmatrix\obsmatrix^\top$ is precomputed once and each step costs one model evaluation plus a cheap $\obsdim\times \obsdim$ solve, identically for all three samplers; with the ensemble-shared inflated covariance (Section~\ref{subsec:approximations}), one Jacobian evaluation and one factorisation per step are shared by the whole ensemble. \emph{Pseudo-time clamping.} The coefficients degenerate at $\tau=0$ ($\beta_0=0$ makes $\gweight_\tau$ singular), so the correction is applied from $\tau=\Delta\tau$ onwards. \emph{Ensemble draws.} An ensemble follows from independent draws (Brownian paths for the SDEs, initial latents for FM); a full trajectory feeds each posterior sample back as $\state^{n}$ and assimilates $\obs^{n+1}$. \emph{Solver choice.} The dynamics are continuous in pseudo-time, so the solver and $M$ may be chosen after training: the guided ODE admits higher-order deterministic solvers, the SDE samplers use Euler--Maruyama.
